\documentclass[12pt,a4paper]{ctexart}
\usepackage[left=2.5cm,right=2.5cm,top=2.3cm,bottom=2.3cm]{geometry}
\usepackage{xcolor}
\usepackage{tcolorbox}
\usepackage{titlesec}
\usepackage{fancyhdr}
\usepackage{enumitem}
\usepackage{tabularx}
\usepackage{amsmath,amssymb}
\usepackage{hyperref}
\tcbuselibrary{skins,breakable}

\definecolor{MainColor}{HTML}{2A6F73}
\definecolor{SecondColor}{HTML}{3CAEA3}
\definecolor{LightMain}{HTML}{EAF7F6}
\definecolor{TitleDark}{HTML}{1F3A40}
\definecolor{TextGray}{HTML}{555555}
\definecolor{BorderGray}{HTML}{CBD5D6}
\definecolor{WarnColor}{HTML}{F2B35D}
\definecolor{ErrorColor}{HTML}{C94C4C}
\definecolor{DefineColor}{HTML}{6C63B7}

\linespread{1.25}
\setlength{\parindent}{2em}
\setlength{\parskip}{0.25em}
\hypersetup{colorlinks=true,linkcolor=MainColor,urlcolor=MainColor}
\pagestyle{fancy}
\fancyhf{}
\fancyfoot[L]{\textcolor{MainColor}{机器学习笔记}}
\fancyfoot[R]{\textcolor{MainColor}{\thepage}}
\renewcommand{\headrulewidth}{0pt}
\renewcommand{\footrulewidth}{0pt}

\newcommand{\topbar}[2]{\noindent{\small\textcolor{MainColor}{#1}}\hfill{\small\bfseries\textcolor{TitleDark}{#2}}\par\vspace{0.35em}{\color{BorderGray}\hrule}\vspace{1.2em}}
\newcommand{\notetitlecard}[6]{
\begin{tcolorbox}[enhanced,colback=white,colframe=BorderGray,boxrule=0.6pt,arc=2mm,left=12pt,right=12pt,top=10pt,bottom=10pt,borderline west={4pt}{0pt}{SecondColor}]
\begin{tabularx}{\linewidth}{X r}
{\small\textcolor{TextGray}{#1}} & {\small\textcolor{TextGray}{#2}} \\[0.9em]
\multicolumn{2}{l}{{\Huge\bfseries\textcolor{MainColor}{#3}}} \\[0.45em]
\multicolumn{2}{l}{{\large\textcolor{SecondColor!90!black}{#4}}} \\[0.45em]
\multicolumn{2}{l}{{\small\textcolor{TextGray}{课程：#5 \qquad 作者：#6}}}
\end{tabularx}\end{tcolorbox}\vspace{1em}}
\newtcolorbox{summarybox}[1][本节摘要]{enhanced,breakable,colback=LightMain,colframe=SecondColor,colbacktitle=SecondColor,coltitle=white,title={#1},fonttitle=\bfseries,boxrule=0.6pt,arc=2mm,left=12pt,right=12pt,top=10pt,bottom=10pt}
\newtcolorbox{definitionbox}[1][定义]{enhanced,breakable,colback=DefineColor!6,colframe=DefineColor,colbacktitle=DefineColor,coltitle=white,title={#1},fonttitle=\bfseries,boxrule=0.6pt,arc=2mm,left=12pt,right=12pt,top=10pt,bottom=10pt}
\newtcolorbox{keybox}[1][重点]{enhanced,breakable,colback=MainColor!5,colframe=MainColor,colbacktitle=MainColor,coltitle=white,title={#1},fonttitle=\bfseries,boxrule=0.6pt,arc=2mm,left=12pt,right=12pt,top=10pt,bottom=10pt}
\newtcolorbox{examplebox}[1][例题]{enhanced,breakable,colback=white,colframe=SecondColor,colbacktitle=SecondColor!85!black,coltitle=white,title={#1},fonttitle=\bfseries,boxrule=0.6pt,arc=2mm,left=12pt,right=12pt,top=10pt,bottom=10pt}
\newtcolorbox{mistakebox}[1][易错点]{enhanced,breakable,colback=ErrorColor!5,colframe=ErrorColor,colbacktitle=ErrorColor,coltitle=white,title={#1},fonttitle=\bfseries,boxrule=0.6pt,arc=2mm,left=12pt,right=12pt,top=10pt,bottom=10pt}
\newtcolorbox{tipbox}[1][提示]{enhanced,breakable,colback=WarnColor!10,colframe=WarnColor,colbacktitle=WarnColor,coltitle=white,title={#1},fonttitle=\bfseries,boxrule=0.6pt,arc=2mm,left=12pt,right=12pt,top=10pt,bottom=10pt}
\titleformat{\section}{\Large\bfseries\color{MainColor}}{\thesection}{0.8em}{}
\titleformat{\subsection}{\large\bfseries\color{TitleDark}}{\thesubsection}{0.8em}{}
\titleformat{\subsubsection}{\normalsize\bfseries\color{SecondColor!80!black}}{\thesubsubsection}{0.8em}{}
\setlist[itemize]{leftmargin=2.2em,itemsep=0.3em,topsep=0.4em}
\setlist[enumerate]{leftmargin=2.2em,itemsep=0.3em,topsep=0.4em}
\newcommand{\important}[1]{\textbf{\textcolor{MainColor}{#1}}}
\newcommand{\warning}[1]{\textbf{\textcolor{ErrorColor}{#1}}}
\newcommand{\concept}[1]{\textbf{\textcolor{DefineColor}{#1}}}

\begin{document}
\topbar{吴恩达机器学习}{学习笔记}
\notetitlecard{Study Notes Series}{2026 Spring}{吴恩达机器学习笔记(9)}{}{机器学习}{C++与草稿纸}

\begin{summarybox}
本周介绍异常检测和推荐系统。异常检测用正常样本建立概率模型，再根据概率密度识别异常；推荐系统则从用户—物品评分中学习用户偏好和物品特征，并用协同过滤与低秩分解产生推荐。
\end{summarybox}

\section{异常检测}
\subsection{问题动机}
\begin{definitionbox}[异常检测设定]
训练集通常主要由正常样本组成，目标是发现概率很低的样本，例如服务器故障、欺诈交易或传感器异常。为每个特征建立概率密度 $p(x_j)$，组合得到
\[
p(x)=\prod_{j=1}^{n}p(x_j).
\]
若 $p(x)<\varepsilon$，就把样本标为异常。
\end{definitionbox}
\subsection{高斯分布模型}
\begin{definitionbox}[单变量高斯分布]
对特征 $x_j$ 估计均值和方差：
\[
\mu_j=\frac1m\sum_{i=1}^{m}x_j^{(i)},\qquad
\sigma_j^2=\frac1m\sum_{i=1}^{m}(x_j^{(i)}-\mu_j)^2.
\]
高斯密度为
\[
p(x_j;\mu_j,\sigma_j^2)=\frac1{\sqrt{2\pi}\sigma_j}
\exp\!\left(-\frac{(x_j-\mu_j)^2}{2\sigma_j^2}\right).
\]
\end{definitionbox}
\begin{keybox}[开发与评价]
用正常训练样本估计参数，在交叉验证集中混入已知正常和异常样本，选择阈值 $\varepsilon$，再在测试集上报告查准率、查全率或 $F_1$。不能只在正常训练集上选择阈值，因为那样无法衡量真正的异常识别能力。
\end{keybox}
\begin{tipbox}[异常检测与监督学习]
异常检测适合异常类别很少、未来异常类型未知的场景；监督学习适合正负样本都较多且未来错误类型与历史相似的场景。若已有大量标注异常样本，直接使用分类器通常更有效。
\end{tipbox}
\subsection{选择特征与多元高斯}
\begin{examplebox}[构造更有辨识力的特征]
若两个特征单独看都近似正常，但它们的组合关系异常，可以增加比值、乘积或差值等特征。特征工程的目标是让异常样本的概率密度明显低于正常样本。
\end{examplebox}
\begin{definitionbox}[多元高斯分布]
当特征相关性很强时，可以直接用协方差矩阵 $\Sigma$ 建立多元高斯模型：
\[
p(x)=\frac1{(2\pi)^{n/2}|\Sigma|^{1/2}}
\exp\!\left(-\frac12(x-\mu)^T\Sigma^{-1}(x-\mu)\right).
\]
它可以学习倾斜的等密度椭圆，但需要更多样本，并要求协方差矩阵可逆或使用稳定的数值处理。
\end{definitionbox}

\section{推荐系统}
\subsection{问题形式化}
\begin{definitionbox}[评分矩阵]
设 $n_u$ 个用户、$n_m$ 个物品，$r(i,j)=1$ 表示用户 $j$ 给物品 $i$ 评分，$y^{(i,j)}$ 为评分值。记用户是否评价过物品的指示量为 $r(i,j)$，目标是预测未评分项目的偏好。
\end{definitionbox}
\subsection{基于内容的推荐}
\begin{keybox}[用户和物品特征]
如果每个物品有特征向量 $x^{(i)}$，可以为用户 $j$ 学习参数 $\theta^{(j)}$，预测
\[
\hat y^{(i,j)}=(\theta^{(j)})^Tx^{(i)}.
\]
这与多个用户各自训练线性回归相似，但需要为用户学习偏好参数。
\end{keybox}
\subsection{协同过滤}
\begin{definitionbox}[同时学习两类参数]
当物品特征未知时，协同过滤同时学习用户参数 $\theta^{(j)}$ 和物品特征 $x^{(i)}$：
\[
\min_{x,\theta}\frac12\sum_{(i,j):r(i,j)=1}
\left((\theta^{(j)})^Tx^{(i)}-y^{(i,j)}\right)^2
\]
加上对两类参数的正则化。用户评分越相似，算法越可能学到相近的用户偏好或物品表示。
\end{definitionbox}
\begin{keybox}[协同过滤的梯度]
对用户参数和物品特征分别求导即可用梯度下降更新：
\[
\theta^{(j)}\leftarrow\theta^{(j)}-\alpha\sum_{i:r(i,j)=1}
\left((\theta^{(j)})^Tx^{(i)}-y^{(i,j)}\right)x^{(i)},
\]
\[
x^{(i)}\leftarrow x^{(i)}-\alpha\sum_{j:r(i,j)=1}
\left((\theta^{(j)})^Tx^{(i)}-y^{(i,j)}\right)\theta^{(j)}.
\]
\end{keybox}
\subsection{低秩矩阵分解与均值归一化}
\begin{tipbox}[向量化预测]
把物品特征组成矩阵 $X$、用户参数组成矩阵 $\Theta$，评分预测矩阵可写成
\[
X\Theta^T.
\]
这是一种低秩近似：用较短的隐向量表示大量用户—物品关系。
\end{tipbox}
\begin{definitionbox}[均值归一化]
对每个物品先计算已知评分的均值 $\mu_i$，训练残差
\[
y^{(i,j)}\leftarrow y^{(i,j)}-\mu_i,
\]
预测后再加回 $\mu_i$。这样可以让初始预测从“该物品的平均受欢迎程度”开始，而不是全部从零开始。
\end{definitionbox}
\begin{mistakebox}[推荐系统的注意事项]
\begin{itemize}
  \item 只在已评分条目上计算代价，不要把未知评分当成 $0$；
  \item 推荐系统中的隐特征通常没有固定语义，不应过度解释单个维度；
  \item 评价指标应模拟真实推荐场景，不能只比较训练集重构误差；
  \item 新用户或新物品缺少历史数据时会遇到冷启动问题，需要额外特征或交互设计。
\end{itemize}
\end{mistakebox}

\section{本周知识脉络}
\begin{keybox}[概率模型与隐变量表示]
异常检测学习“正常数据长什么样”，用低概率阈值发现异常；推荐系统学习“用户和物品如何匹配”，用协同过滤和低秩隐向量预测偏好。两者都需要合理的验证集、稳定的特征表示和面向实际应用的评价指标。
\end{keybox}
\end{document}
